\documentclass{amsart}
% For dealing with references we use the comment environment
\usepackage{verbatim}
\newenvironment{reference}{\comment}{\endcomment}
%\newenvironment{reference}{}{}
\newenvironment{slogan}{\comment}{\endcomment}
\newenvironment{history}{\comment}{\endcomment}

% For commutative diagrams we use Xy-pic
\usepackage[all]{xy}

% We use 2cell for 2-commutative diagrams.
\xyoption{2cell}
\UseAllTwocells

% We use multicol for the list of chapters between chapters
\usepackage{multicol}

% This is generally recommended for better output
\usepackage{lmodern}
\usepackage[T1]{fontenc}

% For cross-file-references

% Package for hypertext links:
\usepackage{hyperref}

% For any local file, say "hello.tex" you want to link to please

% Theorem environments.
%
\theoremstyle{plain}
\newtheorem{theorem}[subsection]{Theorem}
\newtheorem{proposition}[subsection]{Proposition}
\newtheorem{lemma}[subsection]{Lemma}

\theoremstyle{definition}
\newtheorem{definition}[subsection]{Definition}
\newtheorem{example}[subsection]{Example}
\newtheorem{exercise}[subsection]{Exercise}
\newtheorem{situation}[subsection]{Situation}

\theoremstyle{remark}
\newtheorem{remark}[subsection]{Remark}
\newtheorem{remarks}[subsection]{Remarks}

\numberwithin{equation}{subsection}

% Macros
%
\def\lim{\mathop{\mathrm{lim}}\nolimits}
\def\colim{\mathop{\mathrm{colim}}\nolimits}
\def\Spec{\mathop{\mathrm{Spec}}}
\def\Hom{\mathop{\mathrm{Hom}}\nolimits}
\def\Ext{\mathop{\mathrm{Ext}}\nolimits}
\def\SheafHom{\mathop{\mathcal{H}\!\mathit{om}}\nolimits}
\def\SheafExt{\mathop{\mathcal{E}\!\mathit{xt}}\nolimits}
\def\Sch{\mathit{Sch}}
\def\Mor{\mathop{\mathrm{Mor}}\nolimits}
\def\Ob{\mathop{\mathrm{Ob}}\nolimits}
\def\Sh{\mathop{\mathit{Sh}}\nolimits}
\def\NL{\mathop{N\!L}\nolimits}
\def\CH{\mathop{\mathrm{CH}}\nolimits}
\def\proetale{{pro\text{-}\acute{e}tale}}
\def\etale{{\acute{e}tale}}
\def\QCoh{\mathit{QCoh}}
\def\Ker{\mathop{\mathrm{Ker}}}
\def\Im{\mathop{\mathrm{Im}}}
\def\Coker{\mathop{\mathrm{Coker}}}
\def\Coim{\mathop{\mathrm{Coim}}}

% Boxtimes
%
\DeclareMathSymbol{\boxtimes}{\mathbin}{AMSa}{"02}

%
% Macros for moduli stacks/spaces
%
\def\QCohstack{\mathcal{QC}\!\mathit{oh}}
\def\Cohstack{\mathcal{C}\!\mathit{oh}}
\def\Spacesstack{\mathcal{S}\!\mathit{paces}}
\def\Quotfunctor{\mathrm{Quot}}
\def\Hilbfunctor{\mathrm{Hilb}}
\def\Curvesstack{\mathcal{C}\!\mathit{urves}}
\def\Polarizedstack{\mathcal{P}\!\mathit{olarized}}
\def\Complexesstack{\mathcal{C}\!\mathit{omplexes}}
% \Pic is the operator that assigns to X its picard group, usage \Pic(X)
% \Picardstack_{X/B} denotes the Picard stack of X over B
% \Picardfunctor_{X/B} denotes the Picard functor of X over B
\def\Pic{\mathop{\mathrm{Pic}}\nolimits}
\def\Picardstack{\mathcal{P}\!\mathit{ic}}
\def\Picardfunctor{\mathrm{Pic}}
\def\Deformationcategory{\mathcal{D}\!\mathit{ef}}

\setlength{\emergencystretch}{3em}
\begin{document}
\title[Stacks Project, Tag 08YA]{241 --- Decomposition and classification of injective modules over a Noetherian ring}
\maketitle
\noindent Copyright (C) 2005--2025 Johan de Jong. Stacks Project authors. Source: \href{https://stacks.math.columbia.edu/tag/08YA}{Tag 08YA}.

\noindent Editorial difficulty note (not author text). Difficulty H2: The proof constructs a direct-sum decomposition by transfinite recursion, splitting off injective hulls detected by associated primes. The limit-ordinal stages must remain injective and preserve the direct-sum structure; this is a structural classification task beyond routine finite module arguments..

\noindent Editorial provenance note (not author text). History: excerpt from revision \texttt{a0444\allowbreak 6e57e\allowbreak c1fbc\allowbreak 252a8\allowbreak 71afc\allowbreak ec775\allowbreak 2fb28\allowbreak 07b14}, dualizing.tex, lines 645--675. Reference presentation was adapted; original-author context and core support blocks were retained or added. The mathematical wording is unchanged. Licensed under GNU FDL 1.2 or later; no invariant sections or cover texts. See ../licenses/COPYING. Context blocks reproduce author definitions and supporting results; they are not additional counted problems.

\begin{definition}
\label{definition-essential}
Let $\mathcal{A}$ be an abelian category.
\begin{enumerate}
\item An injection $A \subset B$ of $\mathcal{A}$ is {\it essential},
or we say that $B$ is an {\it essential extension of} $A$,
if every nonzero subobject $B' \subset B$ has nonzero intersection with $A$.
\item A surjection $f : A \to B$ of $\mathcal{A}$ is {\it essential}
if for every proper subobject $A' \subset A$ we have $f(A') \not = B$.
\end{enumerate}
\end{definition}

\begin{definition}
\label{definition-injective-hull}
Let $R$ be a ring. An injection $M \to I$ of $R$-modules is said
to be an {\it injective hull} if $I$ is an injective $R$-module and
$M \to I$ is an essential injection.
\end{definition}

\begin{definition}
\label{definition-indecomposable}
An object $X$ of an additive category is called {\it indecomposable}
if it is nonzero and if $X = Y \oplus Z$, then either $Y = 0$ or $Z = 0$.
\end{definition}

\begin{lemma}
\label{lemma-indecomposable-injective-noetherian}
Let $R$ be a Noetherian ring. Let $E$ be an indecomposable injective
$R$-module. Then there exists a prime ideal $\mathfrak p$ of $R$ such that
$E$ is the injective hull of $\kappa(\mathfrak p)$.
\end{lemma}

\begin{proof}
Let $\mathfrak p$ be the prime ideal found in
Lemma \ref{lemma-indecomposable-injective}.
Say $\mathfrak p = (f_1, \ldots, f_r)$.
Pick a nonzero element $x \in \bigcap \Ker(f_i : E \to E)$,
see Lemma \ref{lemma-indecomposable-injective}.
Then $(R_\mathfrak p)x$ is a module isomorphic to $\kappa(\mathfrak p)$
inside $E$. We conclude by Lemma \ref{lemma-indecomposable-injective}.
\end{proof}


\begin{lemma}
\label{lemma-indecomposable-injective}
Let $R$ be a ring. Let $E$ be an indecomposable injective $R$-module.
Then
\begin{enumerate}
\item $E$ is the injective hull of any nonzero submodule of $E$,
\item the intersection of any two nonzero submodules of $E$ is nonzero,
\item $\text{End}_R(E)$ is a noncommutative local ring with maximal
ideal those $\varphi : E \to E$ whose kernel is nonzero, and
\item the set of zerodivisors on $E$ is a prime ideal $\mathfrak p$ of $R$
and $E$ is an injective $R_\mathfrak p$-module.
\end{enumerate}
\end{lemma}

\begin{proof}
Part (1) follows from Lemma \ref{lemma-injective-hull-unique}.
Part (2) follows from part (1) and the definition of injective hulls.

\medskip\noindent
Proof of (3). Set $A = \text{End}_R(E)$ and
$I = \{\varphi \in A \mid \Ker(\varphi) \not = 0\}$.
The statement means that $I$ is a two sided ideal and
that any $\varphi \in A$, $\varphi \not \in I$ is invertible.
Suppose $\varphi$ and $\psi$ are not injective.
Then $\Ker(\varphi) \cap \Ker(\psi)$ is nonzero
by (2). Hence $\varphi + \psi \in I$. It follows that $I$
is a two sided ideal. If $\varphi \in A$, $\varphi \not \in I$,
then $E \cong \varphi(E) \subset E$ is an injective submodule,
hence $E = \varphi(E)$ because $E$ is indecomposable.

\medskip\noindent
Proof of (4). Consider the ring map $R \to A$ and let $\mathfrak p \subset R$
be the inverse image of the maximal ideal $I$. Then it is clear
that $\mathfrak p$ is a prime ideal and that $R \to A$ extends to
$R_\mathfrak p \to A$. Thus $E$ is an $R_\mathfrak p$-module.
It follows from Lemma \ref{lemma-injective-epimorphism} that $E$ is injective
as an $R_\mathfrak p$-module.
\end{proof}


\begin{lemma}
\label{lemma-injective-hull-indecomposable}
Let $\mathfrak p \subset R$ be a prime of a ring $R$.
Let $E$ be the injective hull of $R/\mathfrak p$. Then
\begin{enumerate}
\item $E$ is indecomposable,
\item $E$ is the injective hull of $\kappa(\mathfrak p)$,
\item $E$ is the injective hull of $\kappa(\mathfrak p)$
over the ring $R_\mathfrak p$.
\end{enumerate}
\end{lemma}

\begin{proof}
By Lemma \ref{lemma-essential-extension} the inclusion
$R/\mathfrak p \subset \kappa(\mathfrak p)$ is an essential
extension. Then Lemma \ref{lemma-injective-hull-unique}
shows (2) holds. For $f \in R$, $f \not \in \mathfrak p$
the map $f : \kappa(\mathfrak p) \to \kappa(\mathfrak p)$ is an isomorphism
hence the map $f : E \to E$ is an isomorphism,
see Lemma \ref{lemma-injective-hull-unique}.
Thus $E$ is an $R_\mathfrak p$-module. It is injective
as an $R_\mathfrak p$-module by Lemma \ref{lemma-injective-epimorphism}.
Finally, let $E' \subset E$ be a nonzero injective $R$-submodule.
Then $J = (R/\mathfrak p) \cap E'$ is nonzero. After shrinking $E'$
we may assume that $E'$ is the injective hull of $J$ (see
Lemma \ref{lemma-injective-hull-unique} for example).
Observe that $R/\mathfrak p$ is an essential extension of $J$ for example by
Lemma \ref{lemma-essential-extension}. Hence $E' \to E$
is an isomorphism by Lemma \ref{lemma-injective-hull-unique} part (3).
Hence $E$ is indecomposable.
\end{proof}


\begin{lemma}
\label{lemma-injective-hull}
Let $R$ be a ring. Any $R$-module has an injective hull.
\end{lemma}

\begin{proof}
Let $M$ be an $R$-module. By
More on Algebra, Section \href{https://stacks.math.columbia.edu/tag/01D8}{section injectives modules (Tag 01D8)}
the category of $R$-modules has enough injectives.
Choose an injection $M \to I$ with $I$ an injective $R$-module.
Consider the set $\mathcal{S}$ of submodules $M \subset E \subset I$
such that $E$ is an essential extension of $M$. We order $\mathcal{S}$
by inclusion. If $\{E_\alpha\}$ is a totally ordered subset
of $\mathcal{S}$, then $\bigcup E_\alpha$ is an essential extension of $M$
too (Lemma \ref{lemma-union-essential-extensions}).
Thus we can apply Zorn's lemma and find a maximal element
$E \in \mathcal{S}$. We claim $M \subset E$ is an injective hull, i.e.,
$E$ is an injective $R$-module. This follows from
Lemma \ref{lemma-essential-extensions-in-injective}.
\end{proof}


\begin{lemma}
\label{lemma-injective-hull-unique}
Let $R$ be a ring. Let $M$, $N$ be $R$-modules and let $M \to E$
and $N \to E'$ be injective hulls. Then
\begin{enumerate}
\item for any $R$-module map $\varphi : M \to N$ there exists an
$R$-module map $\psi : E \to E'$ such that
$$
\xymatrix{
M \ar[r] \ar[d]_\varphi & E \ar[d]^\psi \\
N \ar[r] & E'
}
$$
commutes,
\item if $\varphi$ is injective, then $\psi$ is injective,
\item if $\varphi$ is an essential injection, then $\psi$ is an isomorphism,
\item if $\varphi$ is an isomorphism, then $\psi$ is an isomorphism,
\item if $M \to I$ is an embedding of $M$ into an injective $R$-module,
then there is an isomorphism $I \cong E \oplus I'$ compatible with
the embeddings of $M$,
\end{enumerate}
In particular, the injective hull $E$ of $M$ is unique up to isomorphism.
\end{lemma}

\begin{proof}
Part (1) follows from the fact that $E'$ is an injective $R$-module.
Part (2) follows as $\Ker(\psi) \cap M = 0$
and $E$ is an essential extension of $M$.
Assume $\varphi$ is an essential injection. Then
$E \cong \psi(E) \subset E'$ by (2) which implies
$E' = \psi(E) \oplus E''$ because $E$ is injective.
Since $E'$ is an essential extension of
$M$ (Lemma \ref{lemma-essential}) we get $E'' = 0$.
Part (4) is a special case of (3).
Assume $M \to I$ as in (5).
Choose a map $\alpha : E \to I$ extending the map $M \to I$.
Arguing as before we see that $\alpha$ is injective.
Thus as before $\alpha(E)$ splits off from $I$.
This proves (5).
\end{proof}


\begin{lemma}
\label{lemma-essential-extension}
Let $R$ be a ring. Let $M \subset N$ be $R$-modules. The following
are equivalent
\begin{enumerate}
\item $M \subset N$ is an essential extension,
\item for all $x \in N$ nonzero there exists an $f \in R$ such that $fx \in M$
and $fx \not = 0$.
\end{enumerate}
\end{lemma}

\begin{proof}
Assume (1) and let $x \in N$ be a nonzero element. By (1) we have
$Rx \cap M \not = 0$. This implies (2).

\medskip\noindent
Assume (2). Let $N' \subset N$ be a nonzero submodule. Pick $x \in N'$
nonzero. By (2) we can find $f \in R$ with $fx \in M$ and $fx \not = 0$.
Thus $N' \cap M \not = 0$.
\end{proof}


\begin{lemma}
\label{lemma-union-essential-extensions}
Let $R$ be a ring. Let $M$ be an $R$-module. Let $E = \colim E_i$
be a filtered colimit of $R$-modules. Suppose given a compatible
system of essential injections $M \to E_i$ of $R$-modules.
Then $M \to E$ is an essential injection.
\end{lemma}

\begin{proof}
Immediate from the definitions and the fact that filtered
colimits are exact (Algebra, Lemma \href{https://stacks.math.columbia.edu/tag/00DB}{lemma directed colimit exact (Tag 00DB)}).
\end{proof}


\begin{lemma}
\label{lemma-essential-extensions-in-injective}
Let $R$ be a ring. Let $I$ be an injective $R$-module. Let $E \subset I$
be a submodule. The following are equivalent
\begin{enumerate}
\item $E$ is injective, and
\item for all $E \subset E' \subset I$ with $E \subset E'$ essential
we have $E = E'$.
\end{enumerate}
In particular, an $R$-module is injective if and only if every essential
extension is trivial.
\end{lemma}

\begin{proof}
The final assertion follows from the first and the fact that the
category of $R$-modules has enough injectives
(More on Algebra, Section \href{https://stacks.math.columbia.edu/tag/01D8}{section injectives modules (Tag 01D8)}).

\medskip\noindent
Assume (1). Let $E \subset E' \subset I$ as in (2).
Then the map $\text{id}_E : E \to E$ can be extended
to a map $\alpha : E' \to E$. The kernel of $\alpha$ has to be
zero because it intersects $E$ trivially and $E'$ is an essential
extension. Hence $E = E'$.

\medskip\noindent
Assume (2). Let $M \subset N$ be $R$-modules and let $\varphi : M \to E$
be an $R$-module map. In order to prove (1) we have to show that
$\varphi$ extends to a morphism $N \to E$. Consider the set $\mathcal{S}$
of pairs
$(M', \varphi')$ where $M \subset M' \subset N$ and $\varphi' : M' \to E$
is an $R$-module map agreeing with $\varphi$ on $M$. We define an ordering
on $\mathcal{S}$ by the rule $(M', \varphi') \leq (M'', \varphi'')$
if and only if $M' \subset M''$ and $\varphi''|_{M'} = \varphi'$.
It is clear that we can take the maximum of a totally ordered subset
of $\mathcal{S}$. Hence by Zorn's lemma we may assume $(M, \varphi)$
is a maximal element.

\medskip\noindent
Choose an extension $\psi : N \to I$ of $\varphi$ composed
with the inclusion $E \to I$. This is possible as $I$ is injective.
If $\psi(N) \subset E$, then $\psi$ is the desired extension.
If $\psi(N)$ is not contained in $E$, then by (2) the inclusion
$E \subset E + \psi(N)$ is not essential. hence
we can find a nonzero submodule $K \subset E + \psi(N)$ meeting $E$ in $0$.
This means that $M' = \psi^{-1}(E + K)$ strictly contains $M$.
Thus we can extend $\varphi$ to $M'$ using
$$
M' \xrightarrow{\psi|_{M'}} E + K \to (E + K)/K = E
$$
This contradicts the maximality of $(M, \varphi)$.
\end{proof}


\begin{lemma}
\label{lemma-injective-epimorphism}
Let $R \to S$ be an epimorphism of rings. Let $E$ be an $S$-module.
If $E$ is injective as an $R$-module, then $E$ is an injective $S$-module.
\end{lemma}

\begin{proof}
This is true because $\Hom_R(N, E) = \Hom_S(N, E)$ for any $S$-module $N$,
see Algebra, Lemma \href{https://stacks.math.columbia.edu/tag/08YS}{lemma epimorphism modules (Tag 08YS)}.
\end{proof}


\begin{lemma}
\label{lemma-sum-injective-modules}
Let $R$ be a Noetherian ring. A direct sum of injective modules
is injective.
\end{lemma}

\begin{proof}
Let $E_i$ be a family of injective modules parametrized by a set $I$.
Set $E = \bigoplus E_i$. To show that $E$ is injective we use
Injectives, Lemma \href{https://stacks.math.columbia.edu/tag/05NU}{lemma criterion baer (Tag 05NU)}.
Thus let $\varphi : J \to E$ be a module map from an ideal of $R$
into $E$. As $J$ is a finite $R$-module (because $R$ is Noetherian)
we can find finitely many elements $i_1, \ldots, i_r \in I$
such that $\varphi$ maps into $\bigoplus_{j = 1, \ldots, r} E_{i_j}$.
Then we can extend $\varphi$ into $\bigoplus_{j = 1, \ldots, r} E_{i_j}$
using the injectivity of the modules $E_{i_j}$.
\end{proof}


\begin{lemma}
\label{lemma-essential}
Let $\mathcal{A}$ be an abelian category.
\begin{enumerate}
\item If $A \subset B$ and $B \subset C$ are essential extensions, then
$A \subset C$ is an essential extension.
\item If $A \subset B$ is an essential extension and $C \subset B$
is a subobject, then $A \cap C \subset C$ is an essential extension.
\item If $A \to B$ and $B \to C$ are essential surjections, then
$A \to C$ is an essential surjection.
\item Given an essential surjection $f : A \to B$ and a surjection
$A \to C$ with kernel $K$, the morphism $C \to B/f(K)$ is an essential
surjection.
\end{enumerate}
\end{lemma}

\begin{proof}
Omitted.
\end{proof}


\begin{proposition}[Structure of injective modules over Noetherian rings]
\label{proposition-structure-injectives-noetherian}
Let $R$ be a Noetherian ring.
Every injective module is a direct sum of indecomposable injective modules.
Every indecomposable injective module is the injective hull of
the residue field at a prime.
\end{proposition}

\begin{proof}
The second statement is Lemma \ref{lemma-indecomposable-injective-noetherian}.
For the first statement, let $I$ be an injective $R$-module.
We will use transfinite recursion to construct $I_\alpha \subset I$
for ordinals $\alpha$ which are direct sums of indecomposable injective
$R$-modules $E_{\beta + 1}$ for $\beta < \alpha$.
For $\alpha = 0$ we let $I_0 = 0$. Suppose given an ordinal $\alpha$
such that $I_\alpha$ has been constructed. Then $I_\alpha$ is an
injective $R$-module by Lemma \ref{lemma-sum-injective-modules}.
Hence $I \cong I_\alpha \oplus I'$. If $I' = 0$ we are done.
If not, then $I'$ has an associated prime by
Algebra, Lemma \href{https://stacks.math.columbia.edu/tag/0587}{lemma ass zero (Tag 0587)}.
Thus $I'$ contains a copy of $R/\mathfrak p$ for some prime $\mathfrak p$.
Hence $I'$ contains an indecomposable submodule $E$ by
Lemmas \ref{lemma-injective-hull-unique} and
\ref{lemma-injective-hull-indecomposable}. Set
$I_{\alpha + 1} = I_\alpha \oplus E_\alpha$.
If $\alpha$ is a limit ordinal and $I_\beta$ has been constructed
for $\beta < \alpha$, then we set
$I_\alpha = \bigcup_{\beta < \alpha} I_\beta$.
Observe that $I_\alpha = \bigoplus_{\beta < \alpha} E_{\beta + 1}$.
This concludes the proof.
\end{proof}
\end{document}
